
% % -------------------- PROJECTS --------------------
% \section{Personal Projects}
%     \resumeSubHeadingListStart

%     \resumeProjectHeading
%     {\textbf{Yohsin Connect} $|$ \footnotesize\emph{Python, LangChain, FAISS, PyTorch, TensorFlow}}{Ongoing}
%     \resumeItemListStart
%         \resumeItem{Developed a Retrival Augmented Generation (RAG) chatbot to provide instant responses to inquiries from prospective students at Habib University.}
%         \resumeItem{Integrated advanced LLM models using TensorFlow and PyTorch to understand and generate human-like responses.}
%         \resumeItem{Utilized LangChain for streamlined implementation of chat functionalities and FAISS for efficient information retrieval from a large vector database/store of university related documents.}
%         \resumeItem{Ensured seamless operation and scalability of the chatbot by deploying it with Streamlit}
%     \resumeItemListEnd
  
%       \resumeProjectHeading
%         {\textbf{Personal Home Lab and Smart Home Setup} $|$ \footnotesize\emph{YAML, Python, Docker, }}{Ongoing}
%         \resumeItemListStart
%             \resumeItem{Configured diverse services: Jellyfin, Home Assistant, Transmission Daemon with public access}
%             \resumeItem {Managed Docker containers for efficient deployment}
%             % \resumeItem{Deployed SQM QOS on home network to improve bufferbloat}
%         \resumeItemListEnd

    
%         % \resumeProjectHeading
%         % {\textbf{Snake Xenzia Clone} $|$ \footnotesize\emph{ FPGA, VERILOG, XILINX VIVADO}}{September 2021}
%         % \resumeItemListStart
%         %     \resumeItem{Gate level recreation of the infamous snake game, with VGA video out and Cross Matrix USB Keyboard Interfacing for input.}
%         %   \resumeItemListEnd
    
%         \resumeProjectHeading{
%         \href{https://github.com/Talha771/BomberMan-Clone}{\textbf{BomberMan Clone}}
%         $|$ \footnotesize\emph{ C++, SDL2, SPRITE ANIMATIONS, POLYMORPHISM}}{December 2021}
%         \resumeItemListStart
%             \resumeItem{Recreated the famous bomberman game in C++ using SDL2 including sprite animations, and custom game mechanics}
%             \resumeItem{ Implemented OOP design patterns, polymorphism }
        
%           \resumeItemListEnd
          
%       \resumeProjectHeading
%         {\href{https://github.com/Talha771/DSA-Project}{\textbf{Sorting Visualizer}} $|$ \footnotesize\emph{Python,Tkinter,Data Structures}}{August 2022}
%         \resumeItemListStart
%             \resumeItem{Developed a real-time sorting algorithm visualization tool using Tkinter, showcasing sorting algorithms' step-by-step processes to enhance understanding.}
            
%         \resumeItemListEnd

%     \resumeSubHeadingListEnd



% -------------------- PROJECTS --------------------
\section{Projects}
    \resumeSubHeadingListStart

    \resumeProjectHeading
    {\textbf{Yohsin Connect} $|$ \footnotesize\emph{Python, LangChain, FAISS, PyTorch, TensorFlow}}{2024}
    \resumeItemListStart
        \resumeItem{Developed an AI chatbot that instantly responds to prospective students' questions, reducing inquiry response time by 80\%.        .}
        \resumeItem{Integrated TensorFlow and PyTorch LLM models to generate human-like answers, improving response accuracy by 35\%.}
        \resumeItem{Leveraged LangChain and FAISS for optimized information retrieval from large university document databases.}
    \resumeItemListEnd
  
      % \resumeProjectHeading
      %   {\textbf{Personal Home Lab and Smart Home Setup} $|$ \footnotesize\emph{YAML, Python, Docker, }}{Ongoing}
      %   \resumeItemListStart
      %       \resumeItem{Configured diverse services: Jellyfin, Home Assistant, Transmission Daemon with public access}
      %       \resumeItem {Managed Docker containers for efficient deployment}
      %       % \resumeItem{Deployed SQM QOS on home network to improve bufferbloat}
      %   \resumeItemListEnd
    \resumeProjectHeading{
    {\textbf{Custom Unix Shell Implementation}}
      $|$ \footnotesize\emph{C, Unix/Linux Systems Programming, Process Management}}{2024}
      \resumeItemListStart
          \resumeItem{Developed a fully functional Unix shell from scratch, implementing core features like process management, pipe operations, and signal handling}
          \resumeItem{Implemented built-in commands and system calls demonstrating deep understanding of Unix/Linux internals}
      \resumeItemListEnd

    
        % \resumeProjectHeading
        % {\textbf{Snake Xenzia Clone} $|$ \footnotesize\emph{ FPGA, VERILOG, XILINX VIVADO}}{September 2021}
        % \resumeItemListStart
        %     \resumeItem{Gate level recreation of the infamous snake game, with VGA video out and Cross Matrix USB Keyboard Interfacing for input.}
        %   \resumeItemListEnd
    
  \resumeProjectHeading
{\href{https://github.com/Talha771/BomberMan-Clone}{\textbf{BomberMan Clone}} $|$ \footnotesize\emph{C++, SDL2, Sprite Animations, Polymorphism}}{December 2021}
\resumeItemListStart
  \resumeItem{Recreated the classic Bomberman game using C++ and SDL2, incorporating sprite animations and custom game mechanics.}
  \resumeItem{Implemented object-oriented design patterns and polymorphism to enhance code modularity and maintainability.}
  \resumeItem{Developed core gameplay features including player movement, bomb placement, and explosion mechanics.}
\resumeItemListEnd

                    
        \resumeProjectHeading{
        \href{https://github.com/Talha771/BomberMan-Clone}{\textbf{BomberMan Clone}}
        $|$ \footnotesize\emph{ C++, SDL2, SPRITE ANIMATIONS, POLYMORPHISM}}{December 2021}
        \resumeItemListStart
            \resumeItem{Recreated the classic Bomberman game from scratch, adding animations and custom features to make the game more dynamic.}
            \resumeItem{Designed the game using organized coding practices, which improved the code’s flexibility and made it easier to update or expand.}
        
          \resumeItemListEnd
          
    %   \resumeProjectHeading
    %     {\href{https://github.com/Talha771/DSA-Project}{\textbf{Sorting Visualizer}} $|$ \footnotesize\emph{Python,Tkinter,Data Structures}}{August 2022}
    %     \resumeItemListStart
    %         \resumeItem{Developed a real-time sorting algorithm visualization tool using Tkinter, showcasing sorting algorithms' step-by-step processes to enhance understanding.}
            
    %     \resumeItemListEnd

    \resumeSubHeadingListEnd